\section{Language Comparison: Solidity vs ink! vs FunC}

\subsection{Ethereum: Solidity}

Solidity is the dominant smart contract language for EVM-compatible chains. It
is a high-level, curly-bracket language influenced by C++, Python, and
JavaScript.

\begin{lstlisting}[language=Solidity, caption=Chain forwarder in Solidity, label=lst:solidity]
function validateAndForward(
    uint256 chainId,
    bytes calldata data
) external onlyContractA returns (bool) {
    bytes memory transformed = abi.encodePacked(
        "[VALIDATED]", data
    );
    emit MessageValidated(chainId, true, transformed);
    return IContractC(contractC).receiveFinal(
        chainId, transformed
    );
}
\end{lstlisting}

Key traits:
\begin{itemize}
    \item \textbf{Familiar syntax}: Low barrier for students with C/Java background.
    \item \textbf{Mature ecosystem}: Hardhat, Foundry, Remix -- extensive tooling.
    \item \textbf{Interface-based composition}: Cross-contract calls use interface types.
    \item \textbf{Gas-aware design}: Students must optimize storage and computation.
\end{itemize}

\subsection{Polkadot: ink! 5.0}

ink! is a Rust-based eDSL (embedded Domain-Specific Language) for WebAssembly
smart contracts on Substrate-based chains.

\begin{lstlisting}[language=Rust, caption=Chain forwarder in ink! 5.0, label=lst:ink]
fn validate_and_forward(
    &mut self,
    chain_id: u32,
    data: Vec<u8>,
) -> bool {
    let mut transformed = b"[VALIDATED]".to_vec();
    transformed.extend_from_slice(&data);

    self.env().emit_event(MessageValidated {
        chain_id, valid: true,
        transformed_data: transformed.clone(),
    });

    build_call::<DefaultEnvironment>()
        .call(self.contract_c.unwrap())
        .exec_input(
            ExecutionInput::new(
                Selector::new(
                    ink::selector_bytes!("receive_final")
                )
            )
            .push_arg(chain_id)
            .push_arg(transformed),
        )
        .returns::<bool>()
        .invoke()
}
\end{lstlisting}

Key traits:
\begin{itemize}
    \item \textbf{Rust safety}: Ownership, borrowing, and type safety at compile time.
    \item \textbf{Macro-based}: \texttt{\#[ink(storage)]}, \texttt{\#[ink(message)]} annotations.
    \item \textbf{No standard library}: \texttt{no\_std} environment -- no heap by default.
    \item \textbf{Steeper learning curve}: Students need Rust familiarity.
\end{itemize}

\subsection{TON: FunC}

FunC is a low-level language for TON smart contracts. It compiles to Fift
assembler and then to TVM bytecode.

\begin{lstlisting}[language=C, caption=Simple forwarder in FunC, label=lst:func]
() recv_internal(int msg_value, cell in_msg_cell, slice in_msg_body) impure {
    ;; Parse incoming message
    slice ds = in_msg_body;
    int op = ds~load_uint(32);
    int chain_id = ds~load_uint(32);
    slice data = ds~load_dict_slice();

    ;; Validate (non-empty check)
    throw_if(1001, slice_empty?(data));

    ;; Build forward message to ContractB
    cell forward_msg = begin_cell()
        .store_uint(0xba5eba11, 32)
        .store_uint(chain_id, 32)
        .store_slice(data)
        .end_cell();

    ;; Send with bounce flag and remaining balance
    send_raw_message(forward_msg, 64);

    ;; Emit event via storage
    cell event = begin_cell()
        .store_uint(0xcafe, 16)
        .store_uint(chain_id, 32)
        .end_cell();
    set_data(begin_cell().store_ref(event).end_cell());
}
\end{lstlisting}

Key traits:
\begin{itemize}
    \item \textbf{Stack-based}: FunC operates on a stack machine (TVM).
    \item \textbf{Assembly-like}: Manual cell manipulation and gas management.
    \item \textbf{Low-level control}: Full control over message layout and fees.
    \item \textbf{Steepest learning curve}: Most challenging for students.
\end{itemize}

\subsection{Comparative Analysis for Education}

\begin{table}[h]
\centering
\begin{tabular}{|l|c|c|c|}
\hline
\textbf{Criterion} & \textbf{Solidity} & \textbf{ink! 5.0} & \textbf{FunC} \\
\hline
Learning curve & Low & Medium & High \\
Type safety & Low (runtime) & High (compile-time) & Low \\
Tooling maturity & Excellent & Good & Emerging \\
Cross-contract calls & Trivial & Moderate & Complex \\
Debugging & Hardhat/Foundry & Contracts UI & TON CLI \\
Gas estimation & Built-in & Weight system & Manual \\
\hline
\end{tabular}
\caption{Language comparison for smart contract education}
\end{table}

For a blockchain course, the recommended progression is:
\begin{enumerate}
    \item Start with Solidity (familiar syntax, abundant resources).
    \item Move to ink! (safety lessons from Rust, different cross-contract API).
    \item End with FunC (understanding low-level execution, message cells,
          and the actor model at the VM level).
\end{enumerate}

This progression takes students from high-level abstraction to low-level control,
building a complete mental model of how different blockchain VMs approach the
same problem.